\chapter{Podprzestrzenie niezmiennicze i wielomian minimalny}
\label{ch:lin16-niezmiennicze}
\index{podprzestrzeń niezmiennicza@Podprzestrzeń niezmiennicza}
\index{wielomian minimalny@Wielomian minimalny}
\index{przestrzeń cykliczna@Przestrzeń cykliczna}
\index{twierdzenie Cayleya--Hamiltona@Twierdzenie Cayleya--Hamiltona}

\noindent\textbf{Cel rozdziału.}
Diagonalizacja z Rozdziału~\ref{ch:lin14-spektralne}
opisuje operator za pomocą wektorów własnych,
lecz nie każdy operator jest diagonalizowalny.
Badamy więc podprzestrzenie, które operator
zachowuje, oraz wielomiany zerujące operator.
Otrzymamy twierdzenie Cayleya--Hamiltona,
kryterium diagonalizowalności przez wielomian minimalny
i rozkład na podprzestrzenie pierwotne.
Są to narzędzia przygotowujące opis
operatorów niediagonalizowalnych.

W całym rozdziale \(K\) jest dowolnym ciałem,
\(V\) przestrzenią wektorową nad \(K\)
o wymiarze \(n<\infty\),
a \(T\in\operatorname{End}_K(V)\).
Nie zakładamy iloczynu skalarnego.
Dopuszczamy \(n=0\); w tym przypadku
jedynym endomorfizmem \(V\) jest \(I_V=0\).
Dla \(p(t)=\sum_{j=0}^r a_jt^j\in K[t]\)
piszemy \(p(T)=\sum_{j=0}^r a_jT^j\),
gdzie \(T^0=I_V\).

\section{Podprzestrzenie niezmiennicze i operator ilorazowy}
\label{sec:lin16-invariant}

\begin{definicja}[Podprzestrzeń niezmiennicza]
\label{def:lin16-invariant}
Podprzestrzeń \(W\subset V\)
nazywamy \emph{niezmienniczą względem \(T\)},
jeżeli \(T(W)\subset W\).
Wówczas \emph{ograniczenie}
\(T|_W\colon W\to W\) jest endomorfizmem.
Nie wymaga się równości \(T(W)=W\).
\end{definicja}

\begin{przyklad}[Podstawowe podprzestrzenie niezmiennicze]
\label{prz:lin16-basic-invariant}
Podprzestrzenie \(\{0\}\) oraz \(V\)
są niezmiennicze względem każdego \(T\).
Dla \(\lambda\in K\) przestrzeń własna
\(E_\lambda(T)=\ker(T-\lambda I_V)\)
jest niezmiennicza, ponieważ
\(Tv=\lambda v\) dla każdego
\(v\in E_\lambda(T)\).
Dla \(V=K^2\) i
\(T(x,y)=(y,0)\)
prosta \(W=K(1,0)\) jest niezmiennicza,
lecz \(T(W)=\{0\}\ne W\).
\end{przyklad}

\begin{stwierdzenie}[Jądra, obrazy i działania na podprzestrzeniach]
\label{prop:lin16-invariant-kernels}
Niech \(T\in\operatorname{End}_K(V)\).
Dla każdego \(p\in K[t]\)
podprzestrzenie \(\ker p(T)\) oraz
\(\operatorname{im}p(T)\) są niezmiennicze
względem \(T\).
Jeżeli \(W_1,\ldots,W_s\subset V\)
są niezmiennicze względem \(T\),
to ich suma \(W_1+\cdots+W_s\)
i przecięcie \(W_1\cap\cdots\cap W_s\)
również są niezmiennicze.
Ponadto \(p(T)(W_j)\subset W_j\)
dla każdego \(j\).
\end{stwierdzenie}
\begin{proof}
Z definicji wielomianu operatora wynika
\(Tp(T)=p(T)T\), gdyż każda potęga \(T^k\)
przemienia się z \(T\).
Jeśli \(v\in\ker p(T)\), to
\(p(T)(Tv)=Tp(T)v=0\),
więc \(Tv\in\ker p(T)\).
Jeśli \(w=p(T)v\), to
\(Tw=p(T)(Tv)\in\operatorname{im}p(T)\).
Dla wektorów \(w_j\in W_j\)
mamy \(T(\sum_jw_j)=\sum_jTw_j
\in W_1+\cdots+W_s\).
Jeśli \(w\) należy do wszystkich \(W_j\),
to \(Tw\) również należy do każdego z nich.
Wreszcie z \(T(W_j)\subset W_j\)
wynika indukcyjnie \(T^k(W_j)\subset W_j\),
więc analogicznie zachowuje ją \(p(T)\).
\end{proof}

\begin{twierdzenie}[Ograniczenie, iloraz i rozkład blokowy]
\label{thm:lin16-quotient-block}
Niech \(W\subset V\) będzie podprzestrzenią
niezmienniczą względem \(T\).
Wówczas przyporządkowanie
\[
 \overline T\colon V/W\to V/W,\qquad
 \overline T(v+W)=Tv+W
\]
określa endomorfizm liniowy.
Istnieje baza \(V\), w której macierz \(T\)
ma postać
\[
 [T]=\begin{pmatrix}A&B\\0&C\end{pmatrix},
\]
gdzie \(A\) jest macierzą \(T|_W\),
a \(C\) macierzą \(\overline T\).
Dla wielomianów charakterystycznych
zachodzi
\[
 \chi_T(t)=\chi_{T|_W}(t)\,\chi_{\overline T}(t).
\]
\end{twierdzenie}
\begin{proof}
Jeśli \(v+W=v'+W\), to \(v-v'\in W\).
Niezmienniczość daje \(T(v-v')\in W\),
a więc \(Tv+W=Tv'+W\).
Dowodzi to poprawności definicji.
Liniowość wynika z liniowości \(T\).

Niech \((e_1,\ldots,e_k)\)
będzie bazą \(W\), uzupełnioną
do bazy \((e_1,\ldots,e_n)\) przestrzeni \(V\).
Pierwsze \(k\) kolumn macierzy \(T\)
nie mają współrzędnych przy
\(e_{k+1},\ldots,e_n\), gdyż \(T(W)\subset W\).
Stąd dolny lewy blok jest zerowy.
Klasy \(e_{k+1}+W,\ldots,e_n+W\)
stanowią bazę \(V/W\);
działanie \(\overline T\) w tej bazie
jest opisane przez dolny prawy blok \(C\).
Blok \(A\) przedstawia działanie \(T|_W\).

Wyznacznik macierzy blokowo górnotrójkątnej
jest iloczynem wyznaczników bloków diagonalnych,
zatem
\(\det(tI_n-[T])=\det(tI_k-A)\det(tI_{n-k}-C)\).
Gdy \(k=0\) lub \(k=n\), stosujemy
konwencję wyznacznika macierzy pustej równego \(1\).
\end{proof}

\begin{przyklad}[Podprzestrzeń bez niezmienniczego dopełnienia]
\label{prz:lin16-no-invariant-complement}
Niech \(K\) będzie dowolnym ciałem,
\(V=K^2\), a
\[
 T=\begin{pmatrix}\lambda&1\\0&\lambda\end{pmatrix},
 \qquad W=\operatorname{span}\{e_1\},
 \qquad\lambda\in K.
\]
Prosta \(W\) jest niezmiennicza,
a ograniczenie \(T|_W\) i operator
na \(V/W\) są mnożeniem przez \(\lambda\).
W szczególności \(\chi_T(t)=(t-\lambda)^2\).

Nie istnieje jednak
\emph{niezmiennicze} dopełnienie \(W\).
Gdyby \(L\) było takim dopełnieniem,
mielibyśmy \(V=W\oplus L\)
oraz \(\dim L=1\).
Operator na \(L\) działałby
przez pewien skalar \(\mu\in K\),
gdyż \(L\) jest jednowymiarowa.
Skalar \(\mu\) musiałby być wartością własną
\(T\), zatem \(\mu=\lambda\).
Tymczasem \(\ker(T-\lambda I)=W\),
więc \(L\subset W\), wbrew \(V=W\oplus L\).
Pokazuje to, że górny prawy blok \(B\)
nie musi dać się wyzerować przez wybór bazy.
\end{przyklad}

\section{Przestrzenie cykliczne i wielomiany zerujące wektor}
\label{sec:lin16-cyclic}

\begin{definicja}[Przestrzeń cykliczna]
\label{def:lin16-cyclic}
Dla \(v\in V\) definiujemy
\[
 \mathcal K(T,v):=
 \operatorname{span}_K\{v,Tv,T^2v,\ldots\}.
\]
Jest to \emph{przestrzeń cykliczna generowana przez \(v\)}.
Wektor \(v\) nazywamy \emph{cyklicznym dla \(T\)},
jeśli \(\mathcal K(T,v)=V\).
\end{definicja}

\begin{stwierdzenie}[Najmniejsza podprzestrzeń niezmiennicza]
\label{prop:lin16-cyclic-smallest}
Dla każdego \(v\in V\)
przestrzeń \(\mathcal K(T,v)\)
jest najmniejszą podprzestrzenią
niezmienniczą względem \(T\), zawierającą \(v\).
Dla pewnego \(0\le d\le n\) mamy
\[
 \mathcal K(T,v)=
 \operatorname{span}\{v,Tv,\ldots,T^{d-1}v\},
\]
gdzie przy \(d=0\) rozpiętość pustej
rodziny oznacza \(\{0\}\).
Jeżeli \(v\ne0\), można wybrać \(d\ge1\)
tak, że wymienione \(d\) wektorów
tworzy bazę \(\mathcal K(T,v)\).
\end{stwierdzenie}
\begin{proof}
Przestrzeń \(\mathcal K(T,v)\) jest niezmiennicza,
gdyż zastosowanie \(T\) przesuwa każdą
potęgę \(T^jv\) na \(T^{j+1}v\).
Każda podprzestrzeń niezmiennicza
zawierająca \(v\) zawiera indukcyjnie
wszystkie \(T^jv\), a więc ich rozpiętość.

Jeśli \(v=0\), przestrzeń cykliczna jest zerowa
i przyjmujemy \(d=0\).
Jeśli \(v\ne0\), ciąg
\((v,Tv,\ldots,T^nv)\)
zawiera \(n+1\) wektorów w przestrzeni
o wymiarze \(n\), zatem jest liniowo zależny.
Niech \(d\) będzie najmniejszą liczbą dodatnią,
dla której
\[
 T^dv=a_0v+\cdots+a_{d-1}T^{d-1}v
\]
z pewnymi \(a_j\in K\).
Minimalność \(d\) oznacza liniową
niezależność \(v,\ldots,T^{d-1}v\).
Z podanej równości, po kolejnych zastosowaniach
\(T\), wynika indukcyjnie,
że każda potęga \(T^kv\) dla \(k\ge d\)
należy do rozpiętości tych \(d\) wektorów.
Stąd \(d=\dim\mathcal K(T,v)\le n\).
\end{proof}

\begin{definicja}[Wielomian minimalny wektora]
\label{def:lin16-vector-minpoly}
Dla \(v\in V\) \emph{wielomianem minimalnym
wektora \(v\) względem \(T\)}
nazywamy moniczny wielomian \(m_{T,v}\in K[t]\)
najmniejszego stopnia, dla którego
\(m_{T,v}(T)v=0\).
Dla \(v=0\) przyjmujemy \(m_{T,0}=1\).
\end{definicja}

\begin{twierdzenie}[Wielomian minimalny wektora i baza cykliczna]
\label{thm:lin16-vector-minpoly}
Niech \(v\in V\) i
\(d=\dim\mathcal K(T,v)\).
Wielomian \(m_{T,v}\) istnieje, jest jednoznaczny
i ma stopień \(d\).
Dla każdego \(p\in K[t]\) zachodzi
\[
 p(T)v=0\quad\Longleftrightarrow\quad
 m_{T,v}\mid p.
\]
Jeżeli \(v\ne0\), a
\(m_{T,v}(t)=t^d+a_{d-1}t^{d-1}+\cdots+a_0\),
to w bazie \((v,Tv,\ldots,T^{d-1}v)\)
ograniczenie \(T|_{\mathcal K(T,v)}\)
ma dla \(d\ge2\) macierz
\[
 \begin{pmatrix}
 0&0&\cdots&0&-a_0\\
 1&0&\cdots&0&-a_1\\
 0&1&\ddots&0&-a_2\\
 \vdots& &\ddots&\vdots&\vdots\\
 0&0&\cdots&1&-a_{d-1}
 \end{pmatrix}.
\]
Dla \(d=1\) macierz ta ma pojedynczy
element \(-a_0\), to znaczy ma postać
\(\begin{pmatrix}-a_0\end{pmatrix}\).
\end{twierdzenie}
\begin{proof}
Dla \(v=0\) wielomian stały \(1\)
zeruje \(v\), a jego stopień wynosi \(0=d\);
każdy wielomian jest podzielny przez \(1\).
Niech \(v\ne0\).
W dowodzie Stwierdzenia~\ref{prop:lin16-cyclic-smallest}
otrzymaliśmy najmniejsze \(d\ge1\)
oraz współczynniki \(c_0,\ldots,c_{d-1}\)
spełniające \(T^dv=\sum_{j<d}c_jT^jv\).
Wielomian
\(m(t)=t^d-\sum_{j<d}c_jt^j\)
jest moniczny i spełnia \(m(T)v=0\).
Żaden niezerowy wielomian stopnia mniejszego
niż \(d\) nie zeruje \(v\),
bo \(v,\ldots,T^{d-1}v\) są liniowo niezależne.
Wynika stąd także jednoznaczność
monicznego wielomianu najmniejszego stopnia:
różnica dwóch takich wielomianów
miałaby mniejszy stopień i zerowałaby \(v\).

Niech \(p\in K[t]\).
Z dzielenia z resztą mamy
\(p=qm+r\), gdzie \(\deg r<d\) lub \(r=0\).
Po podstawieniu \(T\) i zastosowaniu do \(v\)
uzyskujemy \(p(T)v=r(T)v\).
Jeżeli \(p(T)v=0\), minimalność \(d\)
wymusza \(r=0\), a zatem \(m\mid p\).
Implikacja odwrotna wynika wprost z \(m(T)v=0\).
To pokazuje, że \(m=m_{T,v}\).

Dla \(j<d\) operator \(T\) przesuwa
\(T^{j-1}v\) na następny wektor bazy,
natomiast \(T^dv=-\sum_{j<d}a_jT^jv\).
Odpowiednie kolumny macierzy
mają więc dokładnie wypisaną postać.
\end{proof}

\begin{przyklad}[Wektor cykliczny bloku niediagonalizowalnego]
\label{prz:lin16-jordan-cyclic}
Niech \(V=K^2\) i
\(T=\begin{pmatrix}1&1\\0&1\end{pmatrix}\).
Dla \(v=e_2\) wektory
\(v=(0,1)^{\mathsf T}\) i
\(Tv=(1,1)^{\mathsf T}\) są niezależne,
zatem \(v\) jest cykliczny.
Ponieważ \((T-I)^2=0\), ale
\((T-I)v=e_1\ne0\),
mamy \(m_{T,v}(t)=(t-1)^2\).
Natomiast dla \(w=e_1\)
zachodzi \(Tw=w\), więc
\(\mathcal K(T,w)=Kw\) oraz
\(m_{T,w}(t)=t-1\).
Wielomian minimalny zależy zatem od wektora.
\end{przyklad}

\section{Wielomian minimalny endomorfizmu}
\label{sec:lin16-minpoly}

\begin{definicja}[Wielomian minimalny operatora]
\label{def:lin16-minpoly}
\emph{Wielomianem minimalnym} \(T\)
nazywamy moniczny wielomian
\(m_T\in K[t]\) najmniejszego stopnia,
dla którego \(m_T(T)=0\).
\end{definicja}

\begin{twierdzenie}[Istnienie i własność podzielności]
\label{thm:lin16-minpoly-divides}
Dla każdego endomorfizmu \(T\)
istnieje dokładnie jeden wielomian minimalny \(m_T\).
Dla każdego \(p\in K[t]\)
\[
 p(T)=0\quad\Longleftrightarrow\quad m_T\mid p.
\]
Jeśli \(v_1,\ldots,v_n\) jest bazą \(V\), to
\[
 m_T=\operatorname{lcm}
       (m_{T,v_1},\ldots,m_{T,v_n}),
\]
gdzie najmniejszą wspólną wielokrotność
przyjmujemy moniczną.
Dla \(V=\{0\}\) mamy \(m_T=1\),
a najmniejsza wspólna wielokrotność
pustej rodziny wynosi \(1\).
\end{twierdzenie}
\begin{proof}
Przestrzeń \(\operatorname{End}_K(V)\)
ma wymiar \(n^2\).
Operatory \(I_V,T,T^2,\ldots,T^{n^2}\)
są więc liniowo zależne, gdy \(n\ge1\).
Istnieje niezerowy wielomian \(p\)
spełniający \(p(T)=0\).
Po podzieleniu przez współczynnik wiodący
jest on moniczny.
Wybierzmy spośród takich wielomianów
jeden o najmniejszym stopniu i oznaczmy go \(m_T\).
Jeśli \(m,m'\) są moniczne i mają ten sam
minimalny stopień, ich różnica ma mniejszy stopień
i nadal zeruje \(T\).
Nie może być niezerowa, więc \(m=m'\).

Dla \(p=qm_T+r\) z \(\deg r<\deg m_T\)
lub \(r=0\) otrzymujemy
\(p(T)=r(T)\).
Jeżeli \(p(T)=0\), minimalność
wymusza \(r=0\); odwrotnie wielokrotność
\(m_T\) zeruje \(T\).
To dowodzi podzielności.

Operator \(p(T)\) jest zerowy
dokładnie wtedy, gdy zeruje każdy wektor
wybranej bazy, czyli
gdy \(p(T)v_i=0\) dla wszystkich \(i\).
Na mocy
Twierdzenia~\ref{thm:lin16-vector-minpoly}
jest to równoważne podzielności
\(m_{T,v_i}\mid p\) dla wszystkich \(i\).
A zatem monicznym wielomianem o tej własności,
który dzieli każdy inny, jest ich
najmniejsza wspólna wielokrotność.

Przy \(n=0\) jedyny endomorfizm
jest równocześnie zerowy i identycznościowy,
więc \(1(T)=I_V=0\);
żaden moniczny wielomian nie ma stopnia
ujemnego i definicja daje \(m_T=1\).
\end{proof}

\begin{twierdzenie}[Twierdzenie Cayleya--Hamiltona]
\label{thm:lin16-cayley-hamilton}
Niech \(K\) będzie dowolnym ciałem,
\(V\) skończenie wymiarową przestrzenią
nad \(K\), a \(T\in\operatorname{End}_K(V)\).
Jeżeli \(\chi_T(t)=\det(tI_n-[T])\),
to
\[
 \chi_T(T)=0.
\]
W szczególności \(m_T\mid\chi_T\).
\end{twierdzenie}
\begin{proof}
Dla \(n=0\) mamy \(\chi_T=1\)
oraz \(I_V=0\), więc teza zachodzi.
Niech \(n\ge1\) i wybierzmy bazę \(V\).
Oznaczmy \(A=[T]\) oraz
\(\chi_T(t)=t^n+c_{n-1}t^{n-1}+\cdots+c_0\).
Macierz \(tI_n-A\) ma elementy w
przemiennym pierścieniu \(K[t]\),
dlatego tożsamość dla macierzy dołączonej daje
\[
 (tI_n-A)\operatorname{adj}(tI_n-A)
 =\chi_T(t)I_n.
\]
Każdy element macierzy dołączonej jest
wyznacznikiem pewnej podmacierzy
rozmiaru \(n-1\), toteż jej współczynniki
są wielomianami stopnia co najwyżej \(n-1\).
Zapiszmy
\[
 \operatorname{adj}(tI_n-A)
 =B_0+B_1t+\cdots+B_{n-1}t^{n-1}.
\]
Porównanie współczynników przy potęgach \(t\)
w otrzymanej tożsamości daje
\[
 \begin{aligned}
 -AB_0&=c_0I_n,\\
 B_{j-1}-AB_j&=c_jI_n
       \quad(1\le j\le n-1),\\
 B_{n-1}&=I_n.
 \end{aligned}
\]
Mnożymy pierwszą równość z lewej przez \(I_n\),
równość z indeksem \(j\) przez \(A^j\),
a ostatnią przez \(A^n\).
Po dodaniu wszystkich równości
wyrazy zawierające \(B_j\) skracają się:
\[
 c_0I_n+c_1A+\cdots+c_{n-1}A^{n-1}+A^n=0.
\]
Jest to \(\chi_T(A)=0\), czyli \(\chi_T(T)=0\).
Ważne, że nie podstawialiśmy macierzy
do tożsamości w pierścieniu
nieprzemiennych macierzy wielomianowych:
użyliśmy jedynie porównania współczynników
i mnożenia od lewej.

Ponieważ \(\chi_T(T)=0\), podzielność
\(m_T\mid\chi_T\) wynika z
Twierdzenia~\ref{thm:lin16-minpoly-divides}.
\end{proof}

\begin{wniosek}[Stopień i wartości własne]
\label{cor:lin16-minpoly-spectrum}
Niech \(T\in\operatorname{End}_K(V)\) oraz
\(n=\dim_K V\).
Wówczas \(\deg m_T\le n\), a dla \(n\ge1\)
zachodzi \(\deg m_T\ge1\).
Dla każdego \(\lambda\in K\)
następujące warunki są równoważne:
\[
 \lambda\in\sigma_K(T)
 \quad\Longleftrightarrow\quad
 m_T(\lambda)=0
 \quad\Longleftrightarrow\quad
 \chi_T(\lambda)=0.
\]
\end{wniosek}
\begin{proof}
Z twierdzenia Cayleya--Hamiltona
otrzymujemy \(m_T\mid\chi_T\), co daje
\(\deg m_T\le n\).
Jeśli \(n\ge1\), wielomian stały \(1\)
nie zeruje \(T\), ponieważ \(I_V\ne0\);
zatem \(m_T\) nie jest stały.

Jeżeli \(Tv=\lambda v\) dla \(v\ne0\),
to \(0=m_T(T)v=m_T(\lambda)v\);
stąd \(m_T(\lambda)=0\).
Jeśli \(m_T(\lambda)=0\), podzielność
\(m_T\mid\chi_T\) daje \(\chi_T(\lambda)=0\),
a więc \(\lambda\) jest wartością własną
na mocy
Stwierdzenia~\ref{prop:lin14-charpoly}.
Przy \(n=0\) wszystkie trzy warunki
są fałszywe, a \(m_T=\chi_T=1\).
\end{proof}

\begin{przyklad}[Wielomian minimalny a charakterystyczny]
\label{prz:lin16-minpoly-vs-charpoly}
Niech \(K\) będzie ciałem, w którym
skalary \(1\) i \(2\) są różne.
Dla macierzy diagonalnej
\(D=\operatorname{diag}(1,1,2)\)
mamy
\[
 \chi_D(t)=(t-1)^2(t-2),\qquad
 m_D(t)=(t-1)(t-2).
\]
Istotnie, wielomian zeruje \(D\)
dokładnie wtedy, gdy znika dla obu
wartości własnych \(1\) i \(2\).
Dla macierzy
\[
 A=\begin{pmatrix}
 1&1&0\\0&1&0\\0&0&2
 \end{pmatrix}
\]
wielomian charakterystyczny jest taki sam,
ale
\(m_A(t)=(t-1)^2(t-2)\).
Aby to uzasadnić, zauważmy, że
\(A\) jest sumą bloków
\(J=\begin{pmatrix}1&1\\0&1\end{pmatrix}\)
oraz \((2)\).
Mamy \((J-I)^2=0\), a \(J-I\ne0\),
więc \(m_J=(t-1)^2\);
jednowymiarowy blok ma wielomian minimalny
\(t-2\). Wielomian \(m_A\) jest ich
moniczną najmniejszą wspólną wielokrotnością
na mocy
Twierdzenia~\ref{thm:lin16-minpoly-divides}.
Przykład pokazuje, że samo \(\chi_T\)
nie rozstrzyga struktury operatora.
\end{przyklad}

\section{Rozkład pierwotny podprzestrzeni niezmienniczych}
\label{sec:lin16-primary}

\begin{lemat}[Rozdzielenie względnie pierwszych czynników]
\label{lem:lin16-bezout-kernels}
Niech \(p,q\in K[t]\) będą względnie pierwsze,
a \(T\in\operatorname{End}_K(V)\)
spełniał \(p(T)q(T)=0\).
Wówczas
\[
 V=\ker p(T)\oplus\ker q(T).
\]
Oba składniki są niezmiennicze względem \(T\).
\end{lemat}
\begin{proof}
Ponieważ \(p,q\) są względnie pierwsze,
istnieją \(a,b\in K[t]\) takie, że
\(ap+bq=1\).
Dla \(v\in V\) podstawiamy \(T\):
\[
 v=a(T)p(T)v+b(T)q(T)v.
\]
Pierwszy składnik należy do \(\ker q(T)\),
gdyż wielomiany w \(T\) przemieniają się
i \(p(T)q(T)=0\).
Drugi składnik należy analogicznie
do \(\ker p(T)\).
Ich suma daje więc całe \(V\).
Jeśli \(w\in\ker p(T)\cap\ker q(T)\), to
\(w=a(T)p(T)w+b(T)q(T)w=0\);
suma jest zatem bezpośrednia.
Niezmienniczość wynika ze
Stwierdzenia~\ref{prop:lin16-invariant-kernels}.
\end{proof}

\begin{twierdzenie}[Rozkład pierwotny]
\label{thm:lin16-primary}
Niech \(K\) będzie dowolnym ciałem,
\(V\) skończenie wymiarową przestrzenią
nad \(K\), a \(T\in\operatorname{End}_K(V)\).
Niech
\[
 m_T(t)=\prod_{j=1}^r p_j(t)^{a_j},
\]
gdzie \(p_1,\ldots,p_r\in K[t]\)
są parami różnymi, monicznymi
wielomianami nierozkładalnymi,
a \(a_1,\ldots,a_r\ge1\).
Wtedy
\[
 V=\bigoplus_{j=1}^r
              \ker p_j(T)^{a_j}.
\]
Każdy składnik jest niezmienniczy
względem \(T\).
Dla \(V=\{0\}\) mamy \(r=0\)
i sumę pustą równą \(\{0\}\).
\end{twierdzenie}
\begin{proof}
Przy \(V=\{0\}\) wynik jest bezpośredni.
Niech \(V\ne\{0\}\).
Oznaczmy \(f_j=p_j^{a_j}\).
Ponieważ różne wielomiany nierozkładalne
są względnie pierwsze, również
\(f_i,f_j\) są względnie pierwsze
dla \(i\ne j\).
Wielomian \(m_T=f_1\cdots f_r\)
zeruje \(T\).

Rozkład otrzymamy, rozdzielając kolejne czynniki za pomocą Lematu~\ref{lem:lin16-bezout-kernels}.
Dla \(r=1\) mamy \(f_1(T)=0\), zatem
\(\ker f_1(T)=V\).
Niech \(r\ge2\).
Wielomiany \(f_1\) oraz
\(g=f_2\cdots f_r\) są względnie pierwsze.
Lemat~\ref{lem:lin16-bezout-kernels} daje
\[
 V=\ker f_1(T)\oplus\ker g(T).
\]
Obie podprzestrzenie są niezmiennicze.
Na \(W=\ker g(T)\) operator \(S=T|_W\)
spełnia \(g(S)=0\).
Wielomiany \(f_2,\ldots,f_r\)
są nadal parami względnie pierwsze,
więc możemy ponownie zastosować
Lemat~\ref{lem:lin16-bezout-kernels},
teraz do \(S\) oraz \(f_2\) i
\(f_3\cdots f_r\).
Otrzymujemy
\(W=\ker f_2(S)\oplus
\ker(f_3\cdots f_r)(S)\).
Powtarzając ten sam krok skończenie
wiele razy, uzyskujemy
\[
 W=\bigoplus_{j=2}^r\ker f_j(S).
\]
Ponieważ \(S\) jest ograniczeniem \(T\),
dla \(w\in W\) mamy \(f_j(S)w=f_j(T)w\).
Zatem \(\ker f_j(S)=W\cap\ker f_j(T)\).
Co więcej, jeśli \(w\in\ker f_j(T)\)
dla \(j\ge2\), to \(g(T)w=0\),
bo \(f_j\mid g\), więc \(w\in W\).
Stąd \(\ker f_j(S)=\ker f_j(T)\).
Połączenie otrzymanych sum prostych
kończy indukcję. Niezmienniczość
wszystkich składników wynika
ze Stwierdzenia~\ref{prop:lin16-invariant-kernels}.
\end{proof}

\begin{definicja}[Uogólniona przestrzeń własna]
\label{def:lin16-generalized-eigenspace}
Niech \(K\) będzie ciałem, \(T\in\operatorname{End}_K(V)\),
a \(\lambda\in K\).
Definiujemy \emph{uogólnioną przestrzeń własną}
dla \(\lambda\) przez
\[
 E_\lambda^{\mathrm{gen}}(T)
 :=\bigcup_{k\ge1}\ker(T-\lambda I_V)^k.
\]
Ponieważ ciąg jąder rośnie względem zawierania,
jego suma mnogościowa jest podprzestrzenią \(V\).
Jeśli \(\lambda\notin\sigma_K(T)\),
to \(T-\lambda I_V\) jest odwracalny,
więc \(E_\lambda^{\mathrm{gen}}(T)=\{0\}\).
\end{definicja}

\begin{stwierdzenie}[Stabilizacja uogólnionej przestrzeni własnej]
\label{prop:lin16-gen-stabilize}
Niech \(\lambda\in K\) będzie wartością własną
\(T\in\operatorname{End}_K(V)\)
i niech \((t-\lambda)^a\), \(a\ge1\),
będzie dokładną potęgą tego czynnika
w wielomianie minimalnym \(m_T\).
Wtedy
\[
 \ker(T-\lambda I_V)^a
 =\ker(T-\lambda I_V)^{a+1}
 =\cdots
 =E_\lambda^{\mathrm{gen}}(T).
\]
\end{stwierdzenie}
\begin{proof}
Zapiszmy \(m_T(t)=(t-\lambda)^ag(t)\),
gdzie \(g(\lambda)\ne0\).
Wielomiany \((t-\lambda)^a\) i \(g(t)\)
są względnie pierwsze.
Na mocy Lematu~\ref{lem:lin16-bezout-kernels}
\[
 V=\ker(T-\lambda I_V)^a\oplus\ker g(T).
\]
Na pierwszym składniku \((T-\lambda I)^a=0\).
Na drugim operator \(T-\lambda I\)
jest odwracalny: z tożsamości Bézouta
\(u(t)(t-\lambda)+v(t)g(t)=1\)
wynika, że na \(\ker g(T)\)
operator \(u(T)(T-\lambda I)\) jest identycznością.
Ponieważ czynniki przemieniają się,
\(u(T)\) jest także lewą odwrotnością,
więc \(T-\lambda I\) jest odwracalny
na skończenie wymiarowej \(\ker g(T)\).
W szczególności żadna jego potęga
nie ma tam niezerowego jądra.
Zatem dla każdego \(k\ge a\)
jądro \((T-\lambda I)^k\)
jest dokładnie pierwszym składnikiem,
co dowodzi tezy.
\end{proof}

\begin{wniosek}[Rozkład na uogólnione przestrzenie własne]
\label{cor:lin16-generalized}
Jeżeli wielomian minimalny \(m_T\)
rozkłada się nad \(K\) na czynniki liniowe,
to
\[
 V=\bigoplus_{\lambda\in\sigma_K(T)}
                 E_\lambda^{\mathrm{gen}}(T).
\]
W szczególności zachodzi to nad \(K=\mathbb C\).
\end{wniosek}
\begin{proof}
W rozkładzie pierwotnym
z Twierdzenia~\ref{thm:lin16-primary}
wielomiany \(p_j\) są teraz postaci
\(t-\lambda_j\).
Na mocy
Stwierdzenia~\ref{prop:lin16-gen-stabilize}
odpowiednie jądra potęg są
uogólnionymi przestrzeniami własnymi.
Wartości \(\lambda_j\) są dokładnie
wartościami własnymi, gdyż pierwiastki
\(m_T\) i \(\chi_T\) w \(K\) są te same.
\end{proof}

\begin{przyklad}[Uogólniona przestrzeń własna w wymiarze trzy]
\label{prz:lin16-generalized-example}
Niech \(K=\mathbb C\) oraz
\[
 A=\begin{pmatrix}
 1&1&0\\0&1&0\\0&0&2
 \end{pmatrix}.
\]
Z Przykładu~\ref{prz:lin16-minpoly-vs-charpoly}
mamy \(m_A=(t-1)^2(t-2)\).
Bezpośrednio obliczamy
\[
 \begin{gathered}
 \ker(A-I)=\operatorname{span}\{e_1\},
 \qquad
 \ker(A-I)^2=\operatorname{span}\{e_1,e_2\},\\
 \ker(A-2I)=\operatorname{span}\{e_3\}.
 \end{gathered}
\]
Wobec tego
\(E_1^{\mathrm{gen}}(A)=\operatorname{span}\{e_1,e_2\}\)
i \(E_2^{\mathrm{gen}}(A)=\operatorname{span}\{e_3\}\).
Ich suma jest bezpośrednia i daje \(\mathbb C^3\),
chociaż suma zwykłych przestrzeni własnych
ma tylko wymiar dwa.
\end{przyklad}

\section{Kryterium diagonalizowalności przez wielomian minimalny}
\label{sec:lin16-diagonalization}

\begin{twierdzenie}[Kryterium wielomianu minimalnego]
\label{thm:lin16-minpoly-diagonalizable}
Niech \(K\) będzie dowolnym ciałem,
\(V\) skończenie wymiarową przestrzenią
nad \(K\), a \(T\in\operatorname{End}_K(V)\).
Następujące warunki są równoważne:
\begin{enumerate}[label=(\roman*)]
\item \(T\) jest diagonalizowalny nad \(K\);
\item \(m_T\) jest iloczynem
      parami różnych monicznych czynników liniowych
      nad \(K\).
\end{enumerate}
Przy \(\dim V=0\) warunki zachodzą,
przyjmując iloczyn pusty równy \(1\).
\end{twierdzenie}
\begin{proof}
Załóżmy (i).
W bazie własnej operator \(T\)
ma macierz diagonalną z elementami
\(\lambda_1,\ldots,\lambda_n\in K\).
Niech \(\mu_1,\ldots,\mu_r\)
oznaczają różne wartości na przekątnej.
Wielomian \(p\) zeruje \(T\)
wtedy i tylko wtedy, gdy
\(p(\mu_j)=0\) dla wszystkich \(j\).
Oznacza to podzielność \(t-\mu_j\mid p\)
dla każdego \(j\).
Ponieważ czynniki liniowe są parami
względnie pierwsze, ich iloczyn dzieli \(p\).
Tym samym monicznym wielomianem minimalnym
jest
\(m_T(t)=\prod_{j=1}^r(t-\mu_j)\);
każdy czynnik występuje tylko raz.

Załóżmy (ii) i zapiszmy
\(m_T=\prod_{j=1}^r(t-\mu_j)\)
dla różnych \(\mu_j\in K\).
Wielomiany \(t-\mu_j\) są
parami względnie pierwsze,
a \(m_T(T)=0\).
Twierdzenie~\ref{thm:lin16-primary}
daje
\[
 V=\bigoplus_{j=1}^r
       \ker(T-\mu_jI_V).
\]
Są to zwykłe przestrzenie własne.
Wybierając bazę każdej z nich
i łącząc bazy, otrzymujemy
bazę \(V\) złożoną z wektorów własnych,
czyli (i). Wymiar zero został
uwzględniony przez pustą sumę i pusty iloczyn.
\end{proof}

\begin{uwaga}[Dlaczego wymagamy rozszczepienia nad ciałem]
\label{rem:lin16-split-essential}
Warunek, że \(m_T\) nie ma powtarzających się
czynników nierozkładalnych, nie wystarcza
do diagonalizacji \emph{nad \(K\)},
jeśli czynniki te nie są liniowe.
Na przykład rzeczywista macierz
\[
 R=\begin{pmatrix}0&-1\\1&0\end{pmatrix}
\]
ma \(R^2=-I\), więc
\(m_R(t)=t^2+1\).
Wielomian ten jest nierozkładalny
nad \(\mathbb R\), a zatem nie ma
powtórzonych czynników nierozkładalnych.
Macierz \(R\) nie jest jednak
diagonalizowalna nad \(\mathbb R\),
ponieważ nie ma rzeczywistych
wartości własnych.
Nad \(\mathbb C\) otrzymujemy
\(m_R(t)=(t-i)(t+i)\)
i diagonalizowalność.
\end{uwaga}

\section{Podsumowanie}
\label{sec:lin16-summary}

Podprzestrzeń \(W\subset V\) niezmiennicza
względem \(T\) pozwala rozpatrywać
ograniczenie \(T|_W\) oraz operator
indukowany na \(V/W\).
W bazie dostosowanej do \(W\)
macierz \(T\) jest blokowo górnotrójkątna,
co daje rozkład wielomianu charakterystycznego.
Nie zawsze istnieje niezmiennicze dopełnienie \(W\).

Przestrzeń cykliczna \(\mathcal K(T,v)\)
jest najmniejszą podprzestrzenią niezmienniczą
zawierającą \(v\).
Wielomian minimalny \(m_T\) jest
jedynym monicznym wielomianem
najmniejszego stopnia zerującym \(T\).
Twierdzenie Cayleya--Hamiltona
daje \(m_T\mid\chi_T\), a zatem
\(\deg m_T\le\dim V\).
Wielomian minimalny zachowuje informacje
niewidoczne w samym wielomianie charakterystycznym.

Rozkład pierwotny pozwala przedstawić \(V\)
jako sumę prostą jąder potęg
różnych nierozkładalnych czynników \(m_T\).
Gdy wielomian minimalny rozszczepia się,
otrzymujemy uogólnione przestrzenie własne.
Operator jest diagonalizowalny nad \(K\)
dokładnie wtedy, gdy \(m_T\)
jest iloczynem różnych czynników liniowych
nad tym ciałem.
